% ============================================================
% Laborbuch – 0qln
% SafetySpot Projektdokumentation
% Erstellt auf Basis der Git-Commit-Historie und GitHub-PR-Kommentare
% Formatierung nach IIRu5hII-Vorlage, Styling per safetyspot.sty
% ============================================================

\ssSetDocTitle{Laborbuch}
\ssSetKunde{Bildungsministerium NRW}
\ssSetProjekt{SafetySpot}
\ssSetAutoren{0qln}
\ssSetStand{29.06.2026}
\maketitle
\setcounter{section}{0}

% ============================================================
\section{Grundlage und Projektkontext}
% ============================================================

SafetySpot ist eine mobile Lern-App für Schulen, bei der Schüler interaktiv
Gefahrensituationen aus dem Schulalltag bearbeiten (Chemielabor, Werkraum,
Sportunterricht, Technikraum, Straßenverkehr) und Punkte sammeln.
Lehrer verwalten Klassen und erstellen Szenarien.
Schulen erwerben Lizenzen.

Das Projekt gliedert sich in drei Gradle-Module:
\texttt{ssbackend} (Spring-Boot-REST-API, Java 25),
\texttt{ssmobile} (Android-App, Kotlin/Compose)
und \texttt{sscommon} (geteilte Kotlin-Bibliothek).
Hosting: Microsoft Azure App Service mit Azure SQL für die Produktion;
MariaDB lokal.

Mein persönlicher Schwerpunkt lag auf vier Bereichen:

\begin{enumerate}
  \item \textbf{Infrastruktur und Projektsetup:} Repo-Initialisierung,
        sscommon-Modul, Nix-Dev-Umgebung, CI-Test-Workflow,
        Build-\&-Release-Pipeline für Azure und Android.
  \item \textbf{Backend-API:} Kern-JPA-Entitäten, User-Management-API,
        Image- und Tagging-API, Progress- und Stats-API,
        Admin-Shell und Demo-Daten.
  \item \textbf{Mobile UI Clickdummy:} Vollständiger Compose-Prototyp
        (\textasciitilde{}4\,000 Zeilen Kotlin) mit Auth, Home, Szenarien,
        Gameplay, Ranking und Profil-Screens.
  \item \textbf{Integration:} Backend-Debugging, Szenario-Bilder,
        Merge-Koordination zwischen frontend-functionality und main.
\end{enumerate}

\begin{sshinweis}
  Alle hier genannten Einträge basieren direkt auf Git-Commits und
  GitHub-Pull-Requests im Repository \texttt{safety-spot/SafetySpot}.
  Commit-SHAs und PR-Nummern sind jeweils angegeben.
\end{sshinweis}

% ============================================================
\section{Ausgewertete Unterlagen / Anlagen}
% ============================================================

\begin{ssunterlagetabelle}
Projektvision (docs/v0) &
  Ziele, Stakeholder, Releases und langfristige Vision; Grundlage für
  Scope-Entscheidungen. \\
Informationsmodell (docs/v0) &
  Persistente Datenblöcke: Org/Rollen, Lerninhalte, Fortschritt/Gamification,
  Admin/Lizenz; Grundlage für JPA-Entitäten. \\
Klassendiagramm (docs/v1) &
  Fachliche Klassen- und Beziehungsstruktur; Orientierungsrahmen für
  Backend-Pakete. \\
spec/plan.md &
  Detaillierter Backend-Implementierungsplan: Pakete, DTOs, Security-Config,
  URL-Präfix \texttt{/api/v1}, JWT-Auth-Strategie. \\
spec/issues/ISSUE-001–022 &
  Backend-Issues (001–011) und Mobile-Issues (012–022);
  Grundlage für Aufgabenplanung und PR-Struktur. \\
GitHub-PR-Kommentare &
  z.\,B.\ PR \#44: \glqq{}Waiting for \#31 to complete\grqq{} –
  Koordination mit D0rag3on's Auth-Branch. \\
Spring-Boot-Dokumentation &
  Spring Data JPA, Spring Security (JWT), Spring Shell,
  Spring Boot Actuator. \\
Android/Compose-Dokumentation &
  Jetpack Compose, Navigation-Compose, Material3, Retrofit2. \\
\end{ssunterlagetabelle}

% ============================================================
\section{Technischer Arbeitsstand}
% ============================================================

\begin{ssstandtabelle}
Backend &
  Spring Boot 4.1, Java 25. Ca.\ 93 Java-Dateien / ca.\ 3\,141 Zeilen.
  Pakete: \texttt{auth}, \texttt{user}, \texttt{school}, \texttt{image},
  \texttt{tag}, \texttt{progress}, \texttt{stats}, \texttt{leaderboard},
  \texttt{attempt}, \texttt{dto}, \texttt{enums}, \texttt{exception},
  \texttt{shell}, \texttt{demo}. \\
Mobile Frontend &
  Kotlin/Jetpack Compose. Ca.\ 20 Kotlin-Dateien / ca.\ 4\,108 Zeilen
  (PR \#54). Screens: Auth, Home, Szenarien, Szenario-Detail, Gameplay,
  Ranking, Profil, Settings, Help. \\
Tests &
  18 JUnit-Testklassen in \texttt{ssbackend/src/test/}:
  Controller-, Service- und Repository-Tests für Image, Tag, Progress,
  Stats, School, SchoolClass und User. \\
CI/CD &
  \texttt{test.yml}: Gradle-Tests bei jedem Push/PR.
  \texttt{azure-container-webapp.yml}: Docker-Image → Azure Container Registry
  → Azure App Service.
  \texttt{android-apk-release.yml}: Android APK Build und GitHub Release. \\
Dev-Umgebung &
  Nix devenv (\texttt{devenv.nix}/\texttt{devenv.yaml}/\texttt{devenv.lock})
  mit MariaDB, Java 25 und Gradle; reproduzierbar für alle Teammitglieder. \\
Bildverwaltung &
  Dateibasierte Speicherung in \texttt{ssbackend/data/images/};
  REST-API für Upload/Download; 10 Demo-Szenarien-Bilder enthalten. \\
\end{ssstandtabelle}

% ============================================================
\section{Chronologische Laborbuch-Einträge}
% ============================================================

Klassifikationen orientieren sich an der Laborbuch-Vorlage:
Versuchsaufbau, Entscheidung, Ergebnis, Aufgaben, Problem, Lösung, Notiz.

\begin{sslabortabelle}
26.04.2026 & 0qln & Versuchsaufbau &
  Eintrag 01: Repository \texttt{safety-spot/SafetySpot} initialisiert
  (Commit \texttt{183ac5f}).
  Grundlegende Repo-Struktur: \texttt{.gitignore}, \texttt{.gitattributes},
  \texttt{README.md}; Zeilenendenrichtlinien (\texttt{.gitattributes})
  für plattformübergreifende Zusammenarbeit normalisiert.
  \newline\textit{Ergebnis:} Gemeinsames Repository als Startpunkt für das gesamte Team. \\

09.06.2026 & 0qln & Versuchsaufbau &
  Eintrag 02: Versionsdokumentation eingepflegt – v0-Dokumente (Projektvision,
  Informationsmodell, Entscheidungen, Risiken, Klassendiagramm \texttt{docs/v0})
  und v1-Mockups (\texttt{docs/v1}) dem Repository hinzugefügt
  (Commits \texttt{cf734b4}, \texttt{cf9eabf}).
  \newline\textit{Ergebnis:} Fachliche Grundlage und UX-Mockups im Repo versioniert;
  alle Teammitglieder haben Zugriff ohne separaten Dateiumlauf. \\

09.06.2026 & 0qln & Versuchsaufbau &
  Eintrag 03: \texttt{sscommon}-Kotlin-Modul initiiert sowie
  Copilot-Instruktionen und Projekt-Spezifikation angelegt (PR \#26,
  \texttt{1a90de6}).
  Inhalt: Copilot-Instruktionsdatei \texttt{.github/copilot-instructions.md},
  \texttt{sscommon/}-Gradle-Modul (geteilte Bibliothek),
  \texttt{spec/plan.md} (Backend-Implementierungsplan mit Paketen, DTOs,
  Security-Konzept), ISSUE-001 bis ISSUE-011 (Backend-Aufgaben).
  Außerdem: Dev-Umgebung als Nix-Flake vorbereitet.
  \newline\textit{Ergebnis:} Strukturierte Aufgabenplanung und Copilot-Wissen
  als Grundlage für alle weiteren Backend-Entwicklungsschritte. \\

09.06.2026 & 0qln & Versuchsaufbau &
  Eintrag 04: \texttt{ssmobile}-Android-Projekt initialisiert (PR \#27,
  Commit \texttt{6caaf64}).
  Gradle-Build-Konfiguration, \texttt{AndroidManifest.xml},
  Basis-Theme (Material3), Basis-Navigation und \texttt{MainActivity.kt}.
  Struktur angelegt für spätere UI-Screens.
  \newline\textit{Ergebnis:} Lauffähiges Android-Projekt-Skelett;
  IIRu5hII kann sofort auf der gleichen Basis aufbauen. \\

11.06.2026 & 0qln & Ergebnis &
  Eintrag 05: Projekt-weite Struktur-Überarbeitung (PR \#28, \texttt{9ee221e}).
  Anpassungen an Build-Skripten und Projekt-Struktur nach erstem
  Feedback und paralleler Backend-Entwicklung.
  \newline\textit{Ergebnis:} Konsistente Projektstruktur über alle drei Module. \\

18.06.2026 & 0qln & Aufgaben &
  Eintrag 06: Mobile-Spezifikations-Issues ISSUE-012 bis ISSUE-022 erstellt
  (Commits \texttt{2161ee6}–\texttt{701aa4c}).
  Themen: Foundation/Build-Config, Design-System, Networking/Auth-Infrastruktur,
  Navigation-Shell, Auth-Screen, Home-Screen, Scenarios-Browse,
  Szenario-Gameplay, Ranking-Screen, Profil-Screen, Offline-Cache-Sync.
  \newline\textit{Ergebnis:} Vollständige Spezifikation der mobilen App-Screens
  als Grundlage für IIRu5hII's Frontend-Entwicklung und eigene Clickdummy-Arbeit. \\

21.06.2026 & 0qln & Aufgaben &
  Eintrag 07: Frontend-Planung dokumentiert (PR \#43, \texttt{d01cb56}).
  Entscheidungen zu UI-Architektur, Screen-Übergängen und Datenfluss
  zwischen Frontend-Screens festgehalten.
  \newline\textit{Ergebnis:} Gemeinsames Verständnis für Frontend-Strukturierung
  mit IIRu5hII. \\

22.06.2026 & 0qln & Versuchsaufbau &
  Eintrag 08: Kern-JPA-Entitäten implementiert (PR \#44, \texttt{f5af1c2}).
  Refactoring der Stub-Entitäten von D0rag3on's Basis:
  \texttt{User}, \texttt{School}, \texttt{SchoolClass}, \texttt{Scenario},
  \texttt{Category}, \texttt{Attempt} etc.\ mit Lombok-Annotations
  (\texttt{@Data}, \texttt{@Builder}, \texttt{@Entity}).
  Hibernate konfiguriert (\texttt{spring.jpa.hibernate.ddl-auto=update}).
  Erste Tests zum Laufen gebracht.
  \newline\textit{Notiz:} PR \#44 wartete explizit auf PR \#31 (D0rag3on's Auth):
  \glqq{}Waiting for \#31 to complete\grqq{}. \\

22.06.2026 & 0qln & Ergebnis &
  Eintrag 09: User-Management-API implementiert (PR \#46, \texttt{d8b7aa7}).
  Neue Klassen: \texttt{UserController} (CRUD-Endpoints \texttt{/api/v1/users}),
  \texttt{UserService}/\texttt{UserServiceImpl},
  \texttt{UserRepository} (Spring Data JPA),
  DTOs: \texttt{CreateUserRequest}, \texttt{UpdateUserRequest},
  \texttt{ResetPasswordRequest}, \texttt{UserResponse}.
  Passwörter werden stets gehasht (BCrypt).
  \newline\textit{Ergebnis:} Vollständige CRUD-API für Benutzer; Lehrer können
  Passwörter zurücksetzen. \\

22.06.2026 & 0qln & Ergebnis &
  Eintrag 10: Image- und ImageTag-Entitäten erstellt (PR \#47, \texttt{2c9afa2}).
  \texttt{Image}-Entität mit Metadaten (Beschreibung, Dateiname, Dateipfad),
  \texttt{ImageTag}-Entität (Verknüpfung Bild ↔ Nutzer ↔ Bewertung),
  \texttt{TagValue}-Enum (\texttt{DANGEROUS}/\texttt{SAFE}).
  \newline\textit{Ergebnis:} Datenmodell für kollaboratives Bild-Tagging im Backend. \\

22.06.2026 & 0qln & Ergebnis &
  Eintrag 11: Image-Management-API implementiert (PR \#48, \texttt{7a6cd2f}).
  \texttt{ImageController} (Upload/Download/CRUD \texttt{/api/v1/images}),
  \texttt{ImageService}/\texttt{ImageServiceImpl},
  \texttt{ImageDataService}/\texttt{ImageDataServiceImpl}
  (Datei-basierte Speicherung in \texttt{data/images/}),
  \texttt{ImageRepository}, \texttt{ImageTagRepository}.
  DTOs: \texttt{CreateImageRequest}, \texttt{UpdateImageRequest},
  \texttt{ImageResponse}, \texttt{ImageTagResultResponse}.
  \newline\textit{Ergebnis:} Bilder können hochgeladen, abgerufen und verwaltet werden. \\

22.06.2026 & 0qln & Ergebnis &
  Eintrag 12: Bild-Tagging-API implementiert (PR \#49, \texttt{92ab258}).
  \texttt{TagController} (\texttt{/api/v1/tags}), \texttt{TagService}/
  \texttt{TagServiceImpl}.
  Schüler können Bilder als \texttt{DANGEROUS} oder \texttt{SAFE} markieren;
  Duplikat-Tags werden mit \texttt{DuplicateTagException} abgelehnt.
  DTOs: \texttt{SubmitTagRequest}, \texttt{TagResponse},
  \texttt{ImageStatsResponse}.
  \newline\textit{Ergebnis:} Kern-Spielmechanik (Bild bewerten) im Backend vollständig. \\

22.06.2026 & 0qln & Versuchsaufbau &
  Eintrag 13: Nix Dev-Umgebung eingerichtet (PR \#50, \texttt{3d72b56}).
  \texttt{devenv.nix}: MariaDB, Java 25, Gradle deklarativ definiert;
  \texttt{devenv.yaml} und \texttt{devenv.lock} für Reproduzierbarkeit.
  \newline\textit{Ergebnis:} Einheitliche, reproduzierbare Entwicklungsumgebung
  für alle Teammitglieder ohne manuelle Tool-Installation. \\

22.06.2026 & 0qln & Ergebnis &
  Eintrag 14: Schülerfortschritt- und Lehrer-Stats-API implementiert
  (PR \#51, \texttt{5264857}).
  \texttt{Progress}-Entität (Schüler ↔ Szenario: Score, Attempts, Completed),
  \texttt{ProgressController}/\texttt{ProgressService}/\texttt{ProgressServiceImpl}
  (\texttt{/api/v1/progress}).
  \texttt{StatsController}/\texttt{StatsService}/\texttt{StatsServiceImpl}
  (\texttt{/api/v1/stats}): Klassen-Statistiken pro Lehrer.
  DTOs: \texttt{ProgressEntryResponse}, \texttt{ProgressSummaryResponse},
  \texttt{ClassStatsResponse}, \texttt{StudentStatEntry}.
  \newline\textit{Ergebnis:} Lehrer können Klassenergebnisse einsehen;
  Schüler-Fortschritt wird serverseitig persistiert. \\

22.06.2026 & 0qln & Ergebnis &
  Eintrag 15: Build- \& Release-Pipeline aufgebaut (PR \#52, \texttt{95f8be5}).
  \texttt{Dockerfile} für ssbackend,
  GitHub-Actions-Workflow \texttt{azure-container-webapp.yml}:
  Docker-Build → Push zu Azure Container Registry → Deploy zu Azure App Service.
  \texttt{application-prod.properties} mit Azure-SQL-Datasource-Konfiguration.
  \newline\textit{Ergebnis:} Automatisierter Deployment-Prozess vom Commit
  bis zur Produktionsumgebung auf Azure. \\

22.06.2026 & 0qln & Problem &
  Eintrag 16: Azure-Container-Webapp-Workflow musste direkt nach Anlegen
  zweimal korrigiert werden (Commits \texttt{3ae7804}, \texttt{a9b36ee}).
  Problem: Fehlende Umgebungsvariablen und inkorrekte Workflow-Trigger.
  \newline\textit{Ergebnis:} Workflow läuft stabil nach zwei Korrekturen. \\

23.06.2026 & 0qln & Problem &
  Eintrag 17: Spring-Boot-Tests liefen nicht (Commit-Kontext PR \#53,
  \texttt{3cc8d2a}).
  Ursachen: Fehlende H2-Testdatenbank-Konfiguration, Spring-Security-
  Testkonfiguration nicht gesetzt, Testfixtures unvollständig.
  \newline\textit{Lösung:} H2 In-Memory für Testprofil konfiguriert,
  \texttt{@SpringBootTest}-Annotations ergänzt, 17 Testklassen lauffähig. \\

23.06.2026 & 0qln & Ergebnis &
  Eintrag 18: CI-Test-Workflow erstellt (PR \#56, \texttt{63756d9}).
  \texttt{.github/workflows/test.yml}: Gradle-Tests auf jedem
  Push und Pull-Request automatisch ausgeführt.
  \newline\textit{Ergebnis:} Kontinuierliche Qualitätssicherung;
  Regressions-Erkennung bei jedem Commit. \\

23.06.2026 & 0qln & Ergebnis &
  Eintrag 19: Mobile-UI-Clickdummy vollständig implementiert (PR \#54,
  \texttt{f9ab2db}). Ca.\ 4\,072 Zeilen Kotlin/Compose neu hinzugefügt,
  27 Dateien geändert.
  Screens: \texttt{AuthScreen} (Login/Register), \texttt{HomeScreen}
  (Level, Punkte, Streak, Kategorien), \texttt{ScenariosScreen} (Filter, Suche),
  \texttt{ScenarioDetailScreen}, \texttt{ScenarioPlayScreen} (Gameplay,
  Antwort-Buttons, Feedback), \texttt{RankingScreen} (Scope-Tabs, Top-3),
  \texttt{ProfileScreen} mit Unterseiten.
  Navigation via NavHost/BottomBar, Compose-Theme mit SafetySpot-Farben,
  Mock-Daten für vollständige Offline-Durchläufe.
  \newline\textit{Ergebnis:} Vollständig klickbarer Prototyp zum Demo-Termin bereit. \\

24.06.2026 & 0qln & Ergebnis &
  Eintrag 20: Admin-Shell und Demo-Daten implementiert (PR \#58,
  \texttt{8dc1fe1}).
  \texttt{AdminShellCommands} (Spring Shell): interaktive CLI-Kommandos
  (\texttt{create-school}, \texttt{create-class}, \texttt{create-user},
  \texttt{seed}, \texttt{list-users} etc.).
  \texttt{DemoDataService}: automatischer Seed beim App-Start
  (Schule, Klassen, Schüler, Lehrer, Bilder, Tags).
  \newline\textit{Ergebnis:} Backend kann ohne manuelle SQL-Statements
  mit realistischen Testdaten befüllt werden. \\

24.06.2026 & 0qln & Notiz &
  Eintrag 21: Temporärer Workaround in ssmobile eingeführt (Commit
  \texttt{221ef31}) und am 27.06.2026 wieder rückgängig gemacht
  (Commit \texttt{7a80d24}).
  Hintergrund: Schnelle Umgehung eines Kompatibilitätsproblems zwischen
  Frontend-Functionality-Branch und Backend-Änderungen während
  aktiver Merge-Phase.
  \newline\textit{Ergebnis:} Nur temporär; saubere Lösung nach Merge eingespielt. \\

25.06.2026 & 0qln & Ergebnis &
  Eintrag 22: Demo-Daten korrigiert (PR \#59, \texttt{e766cb6}).
  Fehler im \texttt{DemoDataService}: falsche Beziehungs-IDs und
  inkonsistente Enum-Werte behoben.
  \newline\textit{Ergebnis:} Backend startet mit korrekten und konsistenten
  Testdaten. \\

25.06.2026 & 0qln & Ergebnis &
  Eintrag 23: H2-Konsole für Dev-Profil aktiviert (PR \#60, \texttt{c24bcdf}).
  \texttt{spring.h2.console.enabled=true} in
  \texttt{application-dev.properties}.
  \newline\textit{Ergebnis:} Direkter Datenbankzugriff über Browser während
  Entwicklung erleichtert Debugging. \\

25.06.2026 & 0qln & Ergebnis &
  Eintrag 24: Dateibasierte Bildspeicherung implementiert (PR \#61,
  \texttt{659fa1a}).
  \texttt{ImageDataServiceImpl} schreibt und liest Binärdaten unter
  \texttt{ssbackend/data/images/}.
  Verzeichnis mit gitignore-Ausnahme für Demo-Daten.
  \newline\textit{Ergebnis:} Bilder werden persistent auf dem Dateisystem
  gespeichert und per REST-Endpunkt ausgeliefert. \\

25.06.2026 & 0qln & Aufgaben &
  Eintrag 25: Debugging-Session Backend (Commit \texttt{1d41833}).
  Untersuchung von Fehlern bei der Bild-Upload/Download-Kette zwischen
  ssmobile und ssbackend.
  \newline\textit{Ergebnis:} Probleme identifiziert und in Folgecommits behoben. \\

25.06.2026 & 0qln & Ergebnis &
  Eintrag 26: Szenario-Demo-Bilder hinzugefügt (Commit \texttt{60f19a1}).
  10 PNG-Bilder (\texttt{scenario-00.png} bis \texttt{scenario-09.png})
  unter \texttt{ssbackend/src/main/resources/demo/images/}.
  Diese werden vom \texttt{DemoDataService} beim Seed in die Datenbank geladen.
  \newline\textit{Ergebnis:} Vollständiger Demo-Betrieb mit realistischen
  Bild-Szenarien möglich; IIRu5hII kann Backend-Bilder direkt testen. \\

26.06.2026 & 0qln & Ergebnis &
  Eintrag 27: Projekt-weite Code-Formatierung angewendet (PR \#62,
  \texttt{3bda336}).
  Gradle-Format-Task über alle Java-Dateien in \texttt{ssbackend/}
  ausgeführt; einheitlicher Code-Stil durchgesetzt.
  \newline\textit{Ergebnis:} Diffs in Folge-PRs werden sauberer;
  Review-Aufwand sinkt. \\

27.06.2026 & 0qln & Ergebnis &
  Eintrag 28: Szenario-Bilder aktualisiert (Commit \texttt{4736baf}).
  Bilder-Set überarbeitet und mit dem aktuellen Stand des Backends
  abgeglichen.
  \newline\textit{Ergebnis:} Demo-Bilder im Frontend und Backend konsistent. \\

27.06.2026 & 0qln & Ergebnis &
  Eintrag 29: Android-APK-Release-Workflow erstellt (PR \#64/\#65,
  \texttt{fbaf527}, \texttt{a54247d}, \texttt{485ea21}).
  \texttt{android-apk-release.yml}: automatischer APK-Build bei Git-Tag;
  APK als GitHub-Release-Asset veröffentlicht.
  Zwei Korrekturen nötig (PR \#64 initiale Erstellung, PR \#65 Fixup).
  \newline\textit{Ergebnis:} Vollständige Release-Pipeline für ssmobile;
  APK bei jedem Tag-Push automatisch verfügbar. \\

27.06.2026 & 0qln & Notiz &
  Eintrag 30: Branch \texttt{frontend-functionality} nach umfangreichen
  Merges und Koordination mit IIRu5hII in main zusammengeführt (PR \#63,
  \texttt{0027401}).
  Der PR enthielt Commits beider Autoren (IIRu5hII's Icons/Weitermachen-Logik
  und 0qln's Debugging/Bilder/Workaround-Rücknahme).
  \newline\textit{Ergebnis:} Finaler Stand der mobilen Frontend-Funktionalität
  in main integriert; GitHub-Release \texttt{v0.0.0} erstellt. \\
\end{sslabortabelle}

% ============================================================
\section{Zusammenfassung der Ergebnisse}
% ============================================================

\textbf{Infrastruktur:}
Vollständige CI/CD-Pipeline von Commit bis Produktion:
Test-Workflow, Azure-Deploy-Pipeline und Android-APK-Release.
Reproduzierbare Nix-Dev-Umgebung.

\textbf{Backend-API:}
REST-API mit 8 Hauptbereichen: User, Image, Tag, Progress, Stats,
Leaderboard, School und Auth.
Admin-Shell und Demo-Datenseeding für sofortigen Demo-Betrieb.

\textbf{Mobile Frontend:}
Vollständiger Compose-Clickdummy mit \textasciitilde{}4\,000 Zeilen Kotlin
als Grundlage für IIRu5hII's Backend-Integration.

\textbf{Koordination:}
PR-Kommentare und Merge-Koordination sichergestellt, dass D0rag3on's Auth-Branch
(\#31) vor Core-Entities (\#44) abgeschlossen war.
Gemeinsamer feature-Branch mit IIRu5hII sauber in main gemergt.

\textbf{Lerneffekt:}
Temporäre Workarounds immer sofort mit einem Rücknahme-Commit dokumentieren.
Workflow-Dateien beim ersten Commit testen, nicht raten –
der Azure-Workflow brauchte zwei Korrekturen.

% ============================================================
\section{Auszug aus der Git-Historie von 0qln}
% ============================================================

\begin{ssgittabelle}
26.04.2026 & 183ac5f &
  Initial commit &
  Repo-Initialisierung, README, .gitignore. \\
09.06.2026 & cf734b4 &
  add v0 docs &
  Projektvision, Informationsmodell, Entscheidungen (docs/v0). \\
09.06.2026 & 1a90de6 &
  Initiate clanker setup (\#26) &
  sscommon, Copilot-Instruktionen, spec/plan.md, ISSUE-001–011. \\
09.06.2026 & 826c5bb &
  Initiate ssmobile project (\#27) &
  Android-Projekt-Skeleton mit Gradle, Compose, Navigation. \\
22.06.2026 & f5af1c2 &
  004 – Add core entities (\#44) &
  JPA-Entitäten, Hibernate-Config, Tests. \\
22.06.2026 & d8b7aa7 &
  Implement User Management API (\#46) &
  UserController, Service, Repository, DTOs. \\
22.06.2026 & 2c9afa2 &
  007 image entities (\#47) &
  Image, ImageTag, TagValue-Enum. \\
22.06.2026 & 7a6cd2f &
  008 image management api (\#48) &
  ImageController, DataService (Datei-Speicherung). \\
22.06.2026 & 92ab258 &
  009 – Add image tagging API (\#49) &
  TagController, DANGEROUS/SAFE, DuplicateTagException. \\
22.06.2026 & 3d72b56 &
  Nix dev env (\#50) &
  devenv.nix, MariaDB, Java 25, Gradle. \\
22.06.2026 & 5264857 &
  Add student progress \& teacher stats API (\#51) &
  Progress, Stats, Leaderboard, DTOs. \\
22.06.2026 & 95f8be5 &
  Build- \& Release pipeline (\#52) &
  Dockerfile, Azure-Actions-Workflow, prod.properties. \\
23.06.2026 & 3cc8d2a &
  Get tests to run (\#53) &
  H2-Testprofil, 17 Testklassen lauffähig. \\
23.06.2026 & f9ab2db &
  Mobile UI Clickdummy (\#54) &
  4\,072 Zeilen Kotlin/Compose, 7 Screens. \\
23.06.2026 & 63756d9 &
  Create test.yml (\#56) &
  CI-Workflow: Gradle-Tests bei Push/PR. \\
24.06.2026 & 8dc1fe1 &
  Add admin shell and demo data (\#58) &
  Spring Shell, DemoDataService, Seed-Daten. \\
25.06.2026 & e766cb6 &
  Fix demo data (\#59) &
  Beziehungs-IDs und Enum-Werte korrigiert. \\
25.06.2026 & 659fa1a &
  Image storage (\#61) &
  Dateibasierte Bildspeicherung. \\
25.06.2026 & 60f19a1 &
  Add Scenario images &
  10 Demo-Szenario-Bilder. \\
26.06.2026 & 3bda336 &
  Project wide formatting (\#62) &
  Einheitlicher Code-Stil in ssbackend. \\
27.06.2026 & fbaf527 &
  Create ssmobile release workflow (\#64) &
  Android-APK-Build-\&-Release via GitHub Actions. \\
27.06.2026 & 0027401 &
  Frontend functionality (\#63) &
  Merge frontend-functionality → main; GitHub-Release v0.0.0. \\
\end{ssgittabelle}
